% FILE: content/04_results.tex

\section{Historical Legacy Result Routes and Current Claim Boundaries}

\label{sec:results}

This chapter preserves the result routes inherited from the full Core corpus while preventing those routes from outrunning the current OC Core 1.3.3 proof surface. A legacy theorem route is not automatically a promoted theorem. It becomes release-promoted only where the current claim register, proof sheet, Lean-oriented formalization inventory, finite semantic witness, comparator row, or bounded replay row supports the exact wording and assumptions.


The public reading is therefore explicit. Broad model-vocabulary-level, all-level, scoped, or total-science readings remain historical ambitions unless they are matched by the T133 evidence surface. The routes below are kept because they explain how the model developed and because they give reviewers a map of what has to be checked before a stronger claim may be promoted.


For each route the release asks four questions: what distinction is being attempted, which typed OC construction carries it, which current evidence class can support it, and what condition would reopen it. This turns the older result chapter into a claim-boundary instrument rather than a list of unrestricted theorems.


\subsection{Legacy route 1: dimension-birth route}

The historical text used monotonic dimensionality as a scoped theorem. In the 1.3.3 public surface this is read only as a declared witness route: if a promoted case says that a new axis is born, the proof or finite witness must identify the added axis, the retained lower-level observable, and the condition under which the lower representation loses the relevant verdict.


Promotion rule. The route may appear as a public theorem only when the promoted statement names its assumptions, proof route, executable witness or replay row when relevant, comparator boundary, and falsifier. Without those items it remains a background route in the monograph and cannot be used as evidence for final bounded cross-domain synthesis, full-domain prediction, or unrestricted comparative claim over all modern science.


For legacy route 1, reviewer reopening occurs if the visible public wording is stronger than the current proof sheet, if the finite witness is only self-confirming, if the comparator row absorbs the claimed residual delta, or if a domain example lacks an observable, uncertainty or falsifier.


\subsection{Legacy route 2: axis-availability route}

The historical text said that spontaneous dimension creation is impossible. In the 1.3.3 public surface the promoted claim is narrower: a declared dimensional lift must name an admissible axis or mark the lift as unsupported. This is a model-boundary rule, not an unrestricted physical law.


Promotion rule. The route may appear as a public theorem only when the promoted statement names its assumptions, proof route, executable witness or replay row when relevant, comparator boundary, and falsifier. Without those items it remains a background route in the monograph and cannot be used as evidence for final bounded cross-domain synthesis, full-domain prediction, or unrestricted comparative claim over all modern science.


For legacy route 2, reviewer reopening occurs if the visible public wording is stronger than the current proof sheet, if the finite witness is only self-confirming, if the comparator row absorbs the claimed residual delta, or if a domain example lacks an observable, uncertainty or falsifier.


\subsection{Legacy route 3: death and residue route}

The historical death theorem is retained as a lifecycle route. The promoted 1.3.3 claim is the typed one: liveness, death, residue, and rebirth are separated by the stated status predicates and source-target relations. Any stronger empirical reading requires a domain-specific observable and falsifier.


Promotion rule. The route may appear as a public theorem only when the promoted statement names its assumptions, proof route, executable witness or replay row when relevant, comparator boundary, and falsifier. Without those items it remains a background route in the monograph and cannot be used as evidence for final bounded cross-domain synthesis, full-domain prediction, or unrestricted comparative claim over all modern science.


For legacy route 3, reviewer reopening occurs if the visible public wording is stronger than the current proof sheet, if the finite witness is only self-confirming, if the comparator row absorbs the claimed residual delta, or if a domain example lacks an observable, uncertainty or falsifier.


\subsection{Legacy route 4: embedding-compatibility route}

The historical embedding statements are read as compatibility obligations. A continuum instance may not claim a live realization unless the carrier, realization, boundary, and admissible-space assumptions have been stated. This does not assert a completed ontology of all possible embedding spaces.


Promotion rule. The route may appear as a public theorem only when the promoted statement names its assumptions, proof route, executable witness or replay row when relevant, comparator boundary, and falsifier. Without those items it remains a background route in the monograph and cannot be used as evidence for final bounded cross-domain synthesis, full-domain prediction, or unrestricted comparative claim over all modern science.


For legacy route 4, reviewer reopening occurs if the visible public wording is stronger than the current proof sheet, if the finite witness is only self-confirming, if the comparator row absorbs the claimed residual delta, or if a domain example lacks an observable, uncertainty or falsifier.


\subsection{Legacy route 5: threshold-transition route}

The historical tension-and-threshold language is preserved as a modeling grammar. A promoted theorem or prediction must still define the threshold, the measured or formal quantity compared with it, and the reopening condition. Otherwise the row remains a route, not a finished theorem.


Promotion rule. The route may appear as a public theorem only when the promoted statement names its assumptions, proof route, executable witness or replay row when relevant, comparator boundary, and falsifier. Without those items it remains a background route in the monograph and cannot be used as evidence for final bounded cross-domain synthesis, full-domain prediction, or unrestricted comparative claim over all modern science.


For legacy route 5, reviewer reopening occurs if the visible public wording is stronger than the current proof sheet, if the finite witness is only self-confirming, if the comparator row absorbs the claimed residual delta, or if a domain example lacks an observable, uncertainty or falsifier.


\subsection{Legacy route 6: complexity-growth route}

The historical scoped complexity-growth statement is not promoted as an unrestricted law. In 1.3.3 it is a scoped route: complexity growth may be claimed only for a declared complexity functional, a declared transition, and a proof or replay path showing that the functional changes as stated.


Promotion rule. The route may appear as a public theorem only when the promoted statement names its assumptions, proof route, executable witness or replay row when relevant, comparator boundary, and falsifier. Without those items it remains a background route in the monograph and cannot be used as evidence for final bounded cross-domain synthesis, full-domain prediction, or unrestricted comparative claim over all modern science.


For legacy route 6, reviewer reopening occurs if the visible public wording is stronger than the current proof sheet, if the finite witness is only self-confirming, if the comparator row absorbs the claimed residual delta, or if a domain example lacks an observable, uncertainty or falsifier.


\subsection{Legacy route 7: stabilization route}

The historical no-eternal-stabilization statement is not a public claim about every live system. It is a stress-test route: if a model claims permanent stabilization, it must specify the flows, thresholds, embedding assumptions, and invariant-preservation conditions that make the claim meaningful.


Promotion rule. The route may appear as a public theorem only when the promoted statement names its assumptions, proof route, executable witness or replay row when relevant, comparator boundary, and falsifier. Without those items it remains a background route in the monograph and cannot be used as evidence for final bounded cross-domain synthesis, full-domain prediction, or unrestricted comparative claim over all modern science.


For legacy route 7, reviewer reopening occurs if the visible public wording is stronger than the current proof sheet, if the finite witness is only self-confirming, if the comparator row absorbs the claimed residual delta, or if a domain example lacks an observable, uncertainty or falsifier.


\subsection{Legacy route 8: cycle-support route}

The historical cycle-collapse statements are retained where cycle support is part of the typed model. A promoted claim must distinguish absence of a cycle, failure of a supporting flow, boundary contact, and status change; these are not interchangeable criticism routes.


Promotion rule. The route may appear as a public theorem only when the promoted statement names its assumptions, proof route, executable witness or replay row when relevant, comparator boundary, and falsifier. Without those items it remains a background route in the monograph and cannot be used as evidence for final bounded cross-domain synthesis, full-domain prediction, or unrestricted comparative claim over all modern science.


For legacy route 8, reviewer reopening occurs if the visible public wording is stronger than the current proof sheet, if the finite witness is only self-confirming, if the comparator row absorbs the claimed residual delta, or if a domain example lacks an observable, uncertainty or falsifier.


\subsection{Legacy route 9: expressivity route}

The historical threshold-expressivity route is a useful reviewer test: a representation that cannot express a required distinction cannot carry the corresponding verdict. In 1.3.3 this route is used through theorem-specific assumptions and finite semantic witnesses.


Promotion rule. The route may appear as a public theorem only when the promoted statement names its assumptions, proof route, executable witness or replay row when relevant, comparator boundary, and falsifier. Without those items it remains a background route in the monograph and cannot be used as evidence for final bounded cross-domain synthesis, full-domain prediction, or unrestricted comparative claim over all modern science.


For legacy route 9, reviewer reopening occurs if the visible public wording is stronger than the current proof sheet, if the finite witness is only self-confirming, if the comparator row absorbs the claimed residual delta, or if a domain example lacks an observable, uncertainty or falsifier.


\subsection{Legacy route 10: incompleteness route}

The historical incompleteness-of-embedding-space language is retained only as a boundary discipline. It prevents overclaiming a fixed representation as scoped; it is not itself a proof of scoped OC priority.


Promotion rule. The route may appear as a public theorem only when the promoted statement names its assumptions, proof route, executable witness or replay row when relevant, comparator boundary, and falsifier. Without those items it remains a background route in the monograph and cannot be used as evidence for final bounded cross-domain synthesis, full-domain prediction, or unrestricted comparative claim over all modern science.


For legacy route 10, reviewer reopening occurs if the visible public wording is stronger than the current proof sheet, if the finite witness is only self-confirming, if the comparator row absorbs the claimed residual delta, or if a domain example lacks an observable, uncertainty or falsifier.


\subsection{Legacy route 11: operator route}

The historical operator consequences are promoted only where the current operator semantics, hybrid laws, guard/reset cases, or lifecycle morphisms support the exact public wording.


Promotion rule. The route may appear as a public theorem only when the promoted statement names its assumptions, proof route, executable witness or replay row when relevant, comparator boundary, and falsifier. Without those items it remains a background route in the monograph and cannot be used as evidence for final bounded cross-domain synthesis, full-domain prediction, or unrestricted comparative claim over all modern science.


For legacy route 11, reviewer reopening occurs if the visible public wording is stronger than the current proof sheet, if the finite witness is only self-confirming, if the comparator row absorbs the claimed residual delta, or if a domain example lacks an observable, uncertainty or falsifier.


\subsection{Legacy route 12: cross-level route}

The historical cross-level consequences are read as a K-level witness-taxonomy route. They do not by themselves prove that every adjacent level is ontologically irreducible; irreducibility requires the named retained witness, reduction-failure criterion, demotion criterion, and executable case.


Promotion rule. The route may appear as a public theorem only when the promoted statement names its assumptions, proof route, executable witness or replay row when relevant, comparator boundary, and falsifier. Without those items it remains a background route in the monograph and cannot be used as evidence for final bounded cross-domain synthesis, full-domain prediction, or unrestricted comparative claim over all modern science.


For legacy route 12, reviewer reopening occurs if the visible public wording is stronger than the current proof sheet, if the finite witness is only self-confirming, if the comparator row absorbs the claimed residual delta, or if a domain example lacks an observable, uncertainty or falsifier.


\subsection{Current promoted theorem surface}

The current promoted theorem surface is the T133 proof suite and its supporting artifacts: typed foundation, proof sheets, Lean-oriented formalization inventory, finite semantic checks, and claim-boundary records. The historical routes in this chapter help explain the origin of the questions, but the T133 artifacts carry the current release claims.
